\section{综合讨论：证据链已经推进到哪里？}
\label{sec:discussion}
\subsection{把三个阶段连接起来，而不是把阳性数相加}
整体研究形成了一个相互约束的论证。M0先提供阴性基线，说明正常感觉传播不自动带来当前协议下的跨次记忆；M1给出一个局部慢状态的充分性构造，但其初始传递下降迫使我们检查正常功能。状态干预和旁路随后将“资源历史”与“输出失能”在模型内部区分开。低维代理又说明，即使知道内部资源起作用，单个读出仍可被不同状态规则近似解释，因此需要冻结预测，而不是继续在旧曲线上择优。

前瞻结果支持资源S的有限外推能力，但反弹和hallmark补测共同限制其泛化：平均计数误差达标，不代表全部时序被解释；部分H1/H2/H3/H4(a)/H5证据，也不代表特异性、去习惯化或长期机制已经成立。这里不是五个彼此独立的生物发现，而是\emph{同一候选机制逐步接受更严格约束}。

\begin{table}[!htbp]\centering\small
\caption{需要始终保留的推断边界。证据等级不可仅因报告合并或材料公开而上调。}
\begin{tabularx}{\linewidth}{p{3.1cm}YY}\toprule
已取得的证据 & 可以支持 & 不能直接支持\\\midrule
M0/M1配对比较 & 局部慢资源在本协议下足够 & 全局数学最小性或生物必要性\\
资源救援及快状态对照 & 已构造M1中的历史归因 & 活体JO-CE的分子机制确证\\
训练前后旁路一致 & 神经输出仍可被直接驱动 & 完整动作正常、所有疲劳已排除\\
冻结预测达标 & 选定新条件的粗粒度外推 & 全刺激分布排名或全脑等价替代\\
多项特征方向符合 & 更丰富的习惯化样表型 & 完整动物习惯化或长期学习能力\\
代码/报告公开 & 方法与汇总可检查 & 原始事件已公开或研究已同行评审\\\bottomrule
\end{tabularx}\end{table}

\subsection{Q5：连接组的独特贡献仍是待证明的命题}
全脑M1保留了解剖身份、符号权重、空间路径及循环相互作用，这使细胞级干预成为可定义的任务。代理S只保留一个平均资源状态和简单读出，却已经能够预测若干协议下的aBN1计数。两者并不冲突：在粗粒度输出上可压缩，与在细胞/时序层面具有独特信息，是不同问题。

目前没有运行匹配度数、符号和权重分布的随机图，也没有与局部子回路公平比较，更未检验真实动物的细胞干预。因而既不能说“真实连接组没有用”，也不能说“只有完整真实图才能解释反弹”。结构价值必须通过具体的\emph{新增预测信息}表现出来，而不是通过神经元数量或图谱视觉复杂度获得。

一个有力的下一步，应让候选模型在相同训练数据量、正常功能门槛和评价目标下竞争。例如先用相同计数曲线约束模型，再冻结对某个中间神经元沉默的预测。只有当结构约束在留出干预上提供可复现优势时，才更接近证明连接组信息不可替代。

\subsection{慢状态、参数学习与长期记忆仍不能互换}
M1固定U和恢复常数而改变$x$，这已经足以让当前输出依赖历史；因此“没有运行优化器”不意味着“完全没有记忆”。但历史效应也不自动意味着长期权重学习。2s休息后的H3可以在原参数不变时出现，恰好说明必须区分\emph{从不同残余状态出发}与\emph{学习规则本身改变}。

相同道理，单个规定的2s恢复尺度不能直接等同于神经输出的恢复常数。网络阈值、输入细胞间资源差异和抑制路径都可能改变输出映射。H4(b)的单例80\%提前尤其不应被概括成“高频让突触恢复参数变快”；参数实际上没有改变。现有稀疏点只能约束某些观测阈值时间，不能识别完整动力学谱。

这也说明原研究中的“最小机制”应理解为一种简洁扩展，而不是全局最小性定理。全脑M1按边保存大量资源状态，标量S又使用平均化近似；两者都不是所有可能习惯化系统的唯一形式。

\subsection{哪些首轮局限已经推进，哪些仍未解决}
\begin{table}[!htbp]\centering\small
\caption{逻辑重构后的进展清单。已经执行的内容不再机械复制为“未来工作”，仍未解决的部分则不被成功结果遮蔽。}
\begin{tabularx}{\linewidth}{L{3cm}YY}\toprule
首轮提出的问题 & 已完成的推进 & 保留的缺口\\\midrule
是否只是输出失能？ & 状态救援与强旁路 & 自然输入下低计数；无身体与肌肉测量\\
小模型是否足够？ & 两代理、冻结预测与新协议 & 未比较完整机制家族、局部图与随机图\\
恢复采样太少 & 主补测及一次资格补测 & 时间分辨率有限，未辨识恢复常数\\
多轮/强度未测 & 部分恢复H3与新数据H5 & 真实强度标定、充分恢复后的增强\\
新刺激能否保留/解除响应？ & 两种B及等时/未训练/空白对照 & B的共同读出基线不足，解除未出现\\
真实行为与长期效果 & 尚无相应新增证据 & 输入到动作映射、生理数据与长期协议\\\bottomrule
\end{tabularx}\end{table}

重复回放同一事件阵列有助于排除输入随机变化，但不是自然刺激的全部情形。逐次重新抽样输入、不同初态、多个连接组个体以及生理参数敏感性仍未系统检验。三个种子组也不提供个体动物差异的估计。

\subsection{后续研究的优先级：先补可辨识性，再扩大模型}
\textbf{优先一：建立可比较的跨通路读出。}H7的关键不是随便再选一个B，而是事先有可靠、与A可比较的基线；需要检查细胞标注、输入范围与相同输出，避免根据测试结果不断提高强度直到出现想要的效应。对H8则应另外明确解除刺激的生理依据、时序和一般敏化对照。

\textbf{优先二：以反弹为新预测对象，但不要事后锁定抑制。}在保持平均输入率与初态相同的条件下改变同步性、单细胞分配或时序，再结合细胞或局部结构干预。应在新输出出现前规定反弹幅度、时间或细胞集合指标，并保存竞争预测。旧阶段二称之为“新H4”的研究设想，在本文归入Q5，\emph{尚未执行}，与经典H4(b)补测不是一回事。

\textbf{优先三：比较机制和结构，而非只增加拟合自由度。}可构造神经元适应、慢抑制与多时间尺度资源等受约束候选，统一正常传递门槛、数据预算及干预测试。资源S的成功已有明确机制先验，不能用同一生成模型的预测优势替代活体机制选择。

\textbf{优先四：在有对应数据时连接真实范式。}如果研究问题转为嗅觉抑制或视觉行为，应重新匹配输入编码、时间尺度、输出及腹神经索/身体范围，而不是直接沿用JO-CE的2s常数。糖工程基线也不能直接等同于足部糖的PER习惯化。只有输入、时间和行为读出均可对应，才有资格进行机制鉴别。

以上均为建议，不是本报告的新实验。此前研究按各自预算和停止条件结束；本次重构没有新增仿真或拟合，也不将既有前瞻测试集转为新的训练集。

\section{结论：一个可审计的模型证据链，而不是完整数字动物}
在固定图谱和已测JO-CE协议内，原始LIF模型能传播输入，却没有跨次计数递减。局部可恢复资源补上了部分慢历史效应；随后受控状态干预与旁路进一步将解释限定为该M1内部的资源记忆，而非泛称“全脑学会忽略刺激”。

两个低维代理揭示了输出层面的拟合歧义；运行前冻结的新协议预测则支持资源代理在指定范围内的实用外推，同时保留了其无法精确表示的反弹。最新特征补测把证据从H1/H2/H4(a)扩展到部分恢复H3及人工输入率H5，又明确暴露H4(b)、H6、H7与H8的限制。

\textbf{核心结论是：局部慢状态的条件性充分性、模型内部因果归因和有限前瞻预测可以形成连贯证据，但三者相加仍不等于真实机制确证，也不证明完整连接结构不可替代。}研究应继续用新的、可区分机制与结构的干预，而不是用更多重复的下降曲线替代这一缺口。

\section*{研究透明性与材料范围}\addcontentsline{toc}{section}{研究透明性与材料范围}
\textbf{公开范围。}代码、冻结方案、汇总CSV/JSON、三个历史报告及其排版源码已在项目仓库公开：\url{https://github.com/chenmzh/digital-fly-brain}。原始全脑事件、运行环境、工具缓存及运行日志不随仓库上传；原始事件仅在本机保存，不能把公开汇总的哈希当成已公开的原始数据。

\textbf{版本与时态。}旧报告中“未测”或“未推送”描述当时状态；本稿更新总论，不改写历史PDF、协议或失败记录。各项设计与数据读取顺序保留，不虚构统一的事前预注册。

\textbf{责任与伦理。}本研究仅使用公开来源连接数据进行计算，无新增活体实验，不虚构伦理批准。上游归属与许可按其使用条件遵守。作者、单位、贡献、资助及利益冲突待人工确认。AI助手参与代码、分析和写作，不能替代责任作者审阅；本稿未投稿或经过同行评审。
\FloatBarrier
